\section{Placement in the Grothendieck programme}

\subsection{Cycle conditions and the corrected critical map}
The new detectors provide integral witnesses in the cycle spaces that
enter condition A; the uniform bound supplies their finite-type
parameter space. This is stronger data than the bare assertion that
Lefschetz maps preserve rational algebraic classes. We retain the
correction on cote16, scan050 (archive folio49), in SIGN\_NEXT
SN-C03/C04: on a $2m$-fold $T$, the critical cycle map is $L_T^2$, and
its hyperplane factorization is $j_*L_Wj^*$. The weighted Gysin
condition $L_Tj_*$ cannot be replaced by ordinary lower-critical Gysin.
The historical $C_{\rm G}$ on that page is the algebraicity of the
Lefschetz inverse (modern B), not the modern Kunneth condition.
These distinctions apply to every detector and transfer in this paper;
no repaired handwritten condition is used as a proof of the Hodge
conjecture. The exact modern weak-Lefschetz locator is SGA5, ExposeVII,
Theorem7.1 and its Poincare-dual form; the project correction is the
cited SIGN\_NEXT pair.

The norm bound and the positive Gram matrices in HD-WQ, HD-G05 and
SF-06/07 are the use of condition I on cote16 scan049, and of cote44,
scan002/footer1, Theorem4. They invert a nondegenerate pairing once
geometric detectors exist. The resultant projector in the supplied
Theorem W is the polynomial, Cayley--Hamilton inversion mechanism of
cote16 scans048--049. Our shorter Gaussian and Eisenstein projectors
are explicit integral character-polynomial instances of that mechanism.
The extra cancellation in HD-WQ is a new arithmetic bound on the
multiplier, not a new source of cycles.

\subsection{Schemes over a finitely generated ring over the integers}
All fixed curves, finite covers, group actions, polarizations,
correspondences and effective cycles in a finite portion of the
construction descend to a finitely generated ring over $\Z$ and extend
after shrinking it. This is finite-presentation descent, EGA IV$_3$,
8.8.2; it is not a claim of a smooth model over all of $\operatorname{Spec}\Z$.
For the tame cover constructions invert the group order and the
finitely many level, order and denominator integers involved. The
finite-cover comparison has its standard complex-algebraic placement
in SGA1, ExposeXII, Theorem5.1. Bounded degrees lead to finitely many
cycle parameters as in the accepted HECKE\_DIV R21 and attack R1/D.
The FGA locators used by those proofs are Bourbaki expose221,
Lemme2.4, p.255, and Theoremes3.1--3.2, pp.259--260; one must retain
the bounded-degree argument, rather than cite only a single
fixed-Hilbert-polynomial scheme.

The kernel computations and conductor changes are actual homomorphisms
of abelian schemes after such a base change. Divisor translation,
proper intersection, pullback by a finite etale isogeny and proper
pushforward preserve integral multiplicities in the ways explicitly
proved above. The constants refer to the stated polarization; they
are not assertions independent of an unbounded conductor or lattice.

\subsection{Lambda-rings, K-theory, and the sheaf alternative}
The representation identities for finite groups are identities in
their integral representation rings; exterior powers give the
determinant characters. Their $h^1$ realizations and rational group
projectors are algebraic correspondences. This observation by itself
does not create a middle-dimensional effective cycle. The supplier
and contraction proofs do that additional work.

For a perfect complex the Chern character and exterior operations live
in the K-theoretic intersection framework of SGA6, Expose0, section1,
equations(1.1)--(1.3), p.1, and ExposeV, 2.1--2.2, especially
Lemma2.2.1. Newton and Adams operations have the exact further
locators ExposeV, Appendix7.1, pp.360--361, and Proposition7.5,
pp.362--364. The semiregularity gate of DEG\_L S-14 asks for an actual
perfect complex, a coherent finite translation/tensor action, the
entire equation $e^B\operatorname{ch}(E)=\lambda t+P(\theta)$, and
ordinary negative self-Ext vanishing. A virtual K-class, one Chern
component, or a formal exterior identity does not supply those
hypotheses. Our separate checks keep each hypothesis visible.

\subsection{Etale realizations, weights, and motives}
Smooth proper specialization is SGA4, ExposeXVI, Corollary2.2.
The rational cycle-class compatibility used by the supplied arithmetic
reduction is SGA6, ExposeX, Appendix section7, equation(7.13.10),
pp.578--579. Trace and correspondence operations have the exact
locators SGA5, ExposeIII, section3 and 5.1--5.3. Thus, after the
specified spreading and shrinking, the finite group projectors,
quotient maps and isogeny comparisons act on the lisse sheaves
$R^1f_*\Q_\ell$ and their exterior powers. The Weil summand belongs
to $R^{2n}f_*\Q_\ell(n)$; its untwisted weight is $2n$ and its Tate
twist has weight zero. These are realization statements, not a
deduction of algebraicity from a Frobenius eigenvalue.

In homological motives the rational Weil projector cuts out the
two-dimensional determinant summand. A detecting cycle and its
$K$-endomorphism translate give two independent Tate morphisms into
that summand, with the integral denominators controlled by HD-WQ.
The bounded degree is extra information on representatives and is
not contained in the isomorphism class of the rational motive.
For data over a number field, once these two cycles are defined over
a finite extension, the corresponding good local factor on the
untwisted summand is $(1-q^n u)^{-2}$; before making that extension
there can be a finite Artin action on its Tate twist. This follows
from the fixed algebraic classes and does not estimate any Rallis
value, conductor, Chow degree, or multiplier from an L-function.
No RH or GRH assertion is used. In particular DEG\_L B1's Rallis
point $s=3/2$ remains distinct from a geometric detector estimate.

\subsection{The remaining ghost statement}
The source exclusions D-BI-01--24, D-SPOR, D-SPW, D-G2T,
D-ESC, D-A2, L34, D-HES-F, D-HES-01--06, D-C3-01--03, D-A7-01--03, HD-SPIN, HD-ALT, D-AST, D-EX-01--03, D-UP-01--03,
D-R5-01--04 and HD-E8IF1 concern finite-monodromy coefficient
systems and the dimension of their weight-one period maps. Their
algebraic-geometric objects are the finite covers and rational group
projectors described above; their integral representation-ring and
exterior-power identities belong to the stated lambda-ring framework.
Their semistable graph arguments identify a square-zero monodromy
image in the original coefficient space, and the maximal isotropic
image excludes a nonsplit polarization. The explicit scalar
ramification arguments use finite etale cyclotomic algebras after
inverting the conductor and contrast them with ramification of the
specified coefficient field. These restrictions apply to potential
suppliers for condition A at cote16, scan050, and do not produce
cycles by positivity or extract a bound from a motivic L-function.

SF-US, SF-UG and SF-E811--824 concern actual bundles, coherent sheaves and perfect complexes on a
fourfold carrying genuine real multiplication. The Chern-character
identities are identities in rational cohomology induced by actual
K-classes, in the SGA6 framework specified above. Equivariant
Grothendieck-group additivity, trace and evaluation have the exact
project locator DEG\_L S-21, equations(S21.6)--(S21.7), while the
deformation-input requirements are DEG\_L S-14. The new obstructions
restrict those possible inputs; they neither prove a general negative
algebraicity statement nor identify a ghost. SF-E814 gives an actual
derived example with a different, two-pure character and retains that
scope. Spreading any fixed algebraic construction uses EGA IV$_3$,
8.8.2 as above; the analytic Hermitian--Einstein existence theorem is
not asserted to be a theorem of EGA or SGA. Once an algebraic cycle
exists its etale realization and weights have the stated SGA4/SGA5
placement, but no new cycle or L-function conclusion follows merely
from a Chern--Weil inequality or a self-Ext dimension.

The new bounds close the ghost gate for the Weil planes in the
specified families. Failure of the cyclic or imprimitive construction
is only failure of that source. A counterexample to algebraicity
would give the ghost conclusion only under the exact arithmetic and
CM hypotheses of GHOSTS B02/B03 and HECKE\_DIV P6, with 032 still
conditional as required. We do not use 032 to manufacture the
effective cycles in this paper. Likewise the unbounded core index
in the supplied Theorem E is compatible with our bounded detectors:
fixed, explicitly controlled isogenies in a proof are not a uniform
estimate along an unbounded Hecke transport problem.
